\section{The \v{C}ech obstruction and the two-part contradiction (UC11, UC13)}
\label{sec:obstruction}

This section assembles the obstruction class and states the intended
closing argument. We present the architecture faithfully and then mark,
precisely, which of its parts are established and which are not. The
detailed audit of the two contested inputs is
Section~\ref{sec:proof-state}; here we state where they enter.

\subsection{The minimal counterexample and the trace cover (UC11)}
\label{ssec:trace-cover}

Work in the intersection-closed dialect. Suppose, for contradiction,
that Frankl fails. Among all intersection-closed counterexamples choose
one minimal:

\begin{definition}
\label{def:minimal-cex}
A \emph{minimal counterexample} $\Fstar$ is an intersection-closed
family on $[n]$ with $\bias_x(\Fstar)>0$ for every $x\in[n]$ (no rare
coordinate), minimising $|\Fstar|$ and then $n$.
\end{definition}

\begin{theorem}[Minimality-element theorem, UC11]
\label{thm:minimality-element}
Every \emph{proper} trace $\Fstar|_T$, $T\subsetneq[n]$, has a rare
coordinate: some $x\in T$ with $\bias_x(\Fstar|_T)\le0$.
\end{theorem}

\begin{proof}
$\Fstar|_T$ is a strictly smaller intersection-closed family; if it had
no rare coordinate it would be a smaller counterexample, contradicting
minimality.
\end{proof}

Theorem~\ref{thm:minimality-element} gives a partial section of a
``rare-coordinate'' assignment, defined on every object of the trace
subcategory $\Cn(\Fstar)$ \emph{except} the top object $\Fstar$ itself.
The cohomological machine measures the obstruction to extending it to
$\Fstar$.

\begin{definition}
\label{def:cech-sheaf}
For $x\in[n]$ let $U_x$ be the stratum of proper traces in which $x$ is
rare, $U_x=\{(T,\G)\neq\Fstar:x\in T,\ \bias_x(\G)\le0\}$. By
Theorem~\ref{thm:minimality-element} the strata $\{U_x\}_{x\in[n]}$
cover $\Cn(\Fstar)\setminus\{\Fstar\}$. On $U_x$ put the local Walsh
bias cochain $\widehat b_x=b_x\cdot[\walsh{\{x\}}]$ with
$b_x(\G)=-\bias_x(\G)\ge0$, and on overlaps the \emph{mismatch cochain}
$m_{xy}=\widehat b_x-\widehat b_y$. These data assemble into the
\v{C}ech sheaf $\fobs$ of the obstruction, and into the \v{C}ech
double complex
\[
  \check C^{p,q}(\mathfrak U;\fobs)=\bigoplus_{x_0<\cdots<x_p}
  \fobs^{q}(U_{x_0}\cap\cdots\cap U_{x_p}),
\]
with \v{C}ech differential $\check d$ and cochain differential $d$.
\end{definition}

\begin{lemma}[Cocycle condition, UC11]
\label{lem:cocycle}
$m_{xy}+m_{yz}+m_{zx}=0$ on triple overlaps, so $\{m_{xy}\}$ is a
\v{C}ech $1$-cocycle. Promoted through the edge map of the
\v{C}ech-to-cohomology spectral sequence, it defines a class
\[
  \ob(\Fstar)\in\widetilde H^{n-1}(\Xn;\Q).
\]
\end{lemma}

\begin{remark}
\label{rem:source-level1}
By construction every cochain of $\check C^{p,q}(\mathfrak U;\fobs)$ is
built from level-$1$ Walsh data: $\widehat b_x$ carries the character
$\walsh{\{x\}}$. The \v{C}ech and cochain differentials are additive
and $(\Z/2)^n$-equivariant, so by
Lemma~\ref{lem:equivariant-preserves} the whole double complex --- and
hence $\ob(\Fstar)$ --- stays in level-$1$ isotypes \emph{unless} some
operation multiplies by a Walsh character. The double complex carries
no cup product (it is purely additive). This is the mechanism behind
$Y_2$; see Section~\ref{ssec:y2}.
\end{remark}

\subsection{Fact One: the obstruction is nonzero}
\label{ssec:fact-one}

\begin{theorem}[Non-vanishing, UC11 Lemmas 6.1--6.2]
\label{thm:fact-one}
$\ob(\Fstar)\neq0$ in $\widetilde H^{n-1}(\Xn;\Q)$.
\end{theorem}

\begin{proof}[Proof sketch]
If $\ob(\Fstar)=0$ the mismatch cocycle is a \v{C}ech coboundary, so
the local rare-coordinate data $\{\widehat b_x\}$ glue into a global
section: the partial rare-coordinate assignment of
Theorem~\ref{thm:minimality-element} extends to the top object
$\Fstar$, producing a coordinate $x_\ast$ with
$\bias_{x_\ast}(\Fstar)\le0$. But a minimal counterexample has
$\bias_x(\Fstar)>0$ for every $x$. Hence $\ob(\Fstar)\neq0$.
\end{proof}

\begin{status}
Theorem~\ref{thm:fact-one} is argued cleanly at the level of the
\v{C}ech-cohomology / gluing formalism, modulo
Conjecture~\ref{conj:uc10-1} being available to identify the abutment.
Its load-bearing content is the equivalence ``$\ob(\Fstar)=0\iff$ a
global rare-coordinate section exists'' (UC11 Lemma~6.1). A referee
should note that this equivalence \emph{is} the assertion that the
obstruction theory is faithful; see Remark~\ref{rem:wrapper-risk}.
\end{status}

\subsection{Fact Two: the obstruction is zero --- and where it depends on $Y_1$, $Y_2$}
\label{ssec:fact-two}

The competing claim is that $\ob(\Fstar)$ lands in a part of
$\widetilde H^{n-1}(\Xn;\Q)$ that vanishes. This is a conjunction of
two statements.

\begin{definition}[The two contested inputs]
\label{def:y1-y2}
\leavevmode
\begin{description}[leftmargin=2.6em,style=nextline]
\item[$Y_2$ (landing).] The obstruction class lands in the level-$1$
  Walsh isotypes: $\ob(\Fstar)\in\bigoplus_{x\in[n]}\isotype{\{x\}}{n-1}$
  (UC13~\S2.4.1).
\item[$Y_1$ (lower-Walsh vanishing).] Every lower Walsh isotype
  vanishes in positive degree: $\isotype Sk(\Xn)=0$ for all
  $S\subsetneq[n]$ and $k\ge1$; in particular
  $\isotype{\{x\}}{n-1}=0$ for every $x$ (UC10~\S5.3 / UC13~\S4.5,
  Theorem~4.5.1).
\end{description}
\end{definition}

\begin{proposition}[Fact Two, conditional]
\label{prop:fact-two}
$Y_1\wedge Y_2$ implies $\ob(\Fstar)=0$.
\end{proposition}

\begin{proof}
$Y_2$ places $\ob(\Fstar)$ in $\bigoplus_x\isotype{\{x\}}{n-1}$; $Y_1$
(at $|S|=1$, $k=n-1$) makes each summand zero.
\end{proof}

\begin{status}
Proposition~\ref{prop:fact-two} is a correct two-line deduction; its
\emph{hypotheses} are the issue.
\begin{itemize}[leftmargin=1.4em]
\item $Y_2$ is \emph{mechanism-sound}: the landing rests on
  Lemmas~\ref{lem:equivariant-preserves}--\ref{lem:cup-shifts} and
  Remark~\ref{rem:source-level1}, and this has been machine-checked
  (Section~\ref{ssec:y2}). The earlier competing claim of UC11~\S5.3
  --- a top-isotype landing --- is machine-checked impossible.
\item $Y_1$ is \emph{open}. The proof of $Y_1$ given in UC13~\S4.5 is
  logically invalid (Correction~\ref{corr:trace-exact}); the auxiliary
  lemma UC14~$R2$ advertised as repairing it is false
  (Correction~\ref{corr:r2}). The honest re-derivation reduces $Y_1$ to
  a named open computation (Open residual~\ref{res:y1}). $Y_1$ is
  true-but-unproven.
\end{itemize}
Consequently Fact Two is \emph{not established}: one hypothesis is
mechanism-sound, the other is open.
\end{status}

\subsection{The intended closing argument (UC13 \S7)}
\label{ssec:closing}

The two facts are intended to collide. The closing argument of UC13~\S7
is the following seven steps.

\begin{enumerate}[label=\textup{(C\arabic*)},leftmargin=3.1em]
\item Assume Frankl fails; fix a minimal counterexample $\Fstar$
  (Definition~\ref{def:minimal-cex}).
\item Conjecture~\ref{conj:uc10-1} (UC10.1): the lower isotypes
  $\isotype{\{x\}}{n-1}$ vanish.
\item Lemma~\ref{lem:cocycle}: $\ob(\Fstar)$ is defined in
  $\widetilde H^{n-1}(\Xn;\Q)$.
\item Theorem~\ref{thm:fact-one} (Fact One): $\ob(\Fstar)\neq0$.
\item $Y_2$: $\ob(\Fstar)\in\bigoplus_x\isotype{\{x\}}{n-1}$.
\item Steps (C2),(C5): $\ob(\Fstar)=0$ (this is Fact Two,
  Proposition~\ref{prop:fact-two}).
\item Steps (C4),(C6) contradict; hence no minimal counterexample;
  hence Frankl.
\end{enumerate}

This is a clean two-part contradiction: \emph{must-be-nonzero} (from
minimality) versus \emph{must-be-zero} (from the isotype landing plus
the lower-Walsh vanishing). It is the program's intended proof of
Frankl.

\begin{status}
\label{status:closing}
The architecture (C1)--(C7) is coherent, and steps (C1), (C3), (C4),
(C5) are sound or mechanism-sound. The argument does \emph{not} close,
because step (C2) --- equivalently the $Y_1$ half of (C6) --- depends on
Conjecture~\ref{conj:uc10-1}, whose lower-Walsh half $Y_1$ is open
(Section~\ref{ssec:y1}). The UC13 note carried an internal banner
``\GREEN{Frankl closes unconditionally}''; that banner is
\textbf{withdrawn} by the audit of Section~\ref{sec:proof-state}. The
honest status of the closing argument is: a valid proof \emph{modulo}
$Y_1$, with $Y_1$ a precisely isolated open problem.
\end{status}

\begin{remark}[A structural caution for the referee]
\label{rem:wrapper-risk}
The vanishing side of the contradiction uses only the shape of the
\v{C}ech double complex of $\fobs$, never the counterexample
hypothesis $\bias_x(\Fstar)>0$. So, \emph{if} the vanishing argument is
correct, $\ob(\F)$ is identically zero for \emph{every} family, and all
the Frankl-specific content is carried by Fact One --- specifically by
the UC11 Lemma~6.1 equivalence ``$\ob=0\iff$ a global rare-coordinate
section exists''. That equivalence is itself close to a restatement of
Frankl. This is the classic failure mode of cohomological-obstruction
attacks on hard combinatorial problems: the obstruction is
\emph{defined}, \emph{asserted} faithful, and \emph{asserted} to lie in
a zero group; slack in any of the three voids the proof. A referee
should weigh whether UC11 Lemma~6.1 is pinned down to referee standard
independently of the cohomological wrapper. The present paper does not
claim it is.
\end{remark}
